\documentclass[12pt,a4paper]{article}
\usepackage[utf8]{inputenc}
\usepackage[T1]{fontenc}
\usepackage[brazilian]{babel}
\usepackage{geometry}
\usepackage{listings}
\usepackage{xcolor}
\usepackage{hyperref}
\usepackage{fancyhdr}

\geometry{margin=2.5cm}

\lstset{
  backgroundcolor=\color{gray!10},
  basicstyle=\ttfamily\small,
  breaklines=true,
  frame=single,
  rulecolor=\color{gray!40},
  xleftmargin=4pt,
  xrightmargin=4pt,
  aboveskip=8pt,
  belowskip=8pt,
}

\pagestyle{fancy}
\fancyhf{}
\rhead{Sistemas Operacionais}
\lhead{IFSC -- Campus São José}
\rfoot{\thepage}

% ── Preencha aqui ──
\newcommand{\aluno}{Victor E. L. Guerra}
\newcommand{\dataPratica}{27/03/2026}

\begin{document}

\begin{titlepage}
\centering
\vspace*{2cm}
{\Large\textbf{Instituto Federal de Santa Catarina}\\Campus São José\par}
\vspace{0.5cm}
{\large Curso de Engenharia de Telecomunicações\par}
\vspace{3cm}
{\LARGE\textbf{Prática 2 -- Gestão de Tarefas}\par}
\vspace{0.5cm}
{\large Questionário de Entrega\par}
\vspace{3cm}
{\large\textbf{Aluno:} \aluno\par}
\vspace{0.5cm}
{\large\textbf{Professor:} Rogério Pereira Junior\par}
\vspace{3cm}
{\large São José, \dataPratica\par}
\end{titlepage}

\section*{Questionário de Entrega}

\subsection*{1. Qual a diferença entre os estados Ready (Pronto) e Running (Executando) observada durante o uso do comando top?}

No estado \textbf{Ready (Pronto)}, o processo está carregado na memória e apto a executar, porém aguarda que o escalonador lhe atribua a CPU. No estado \textbf{Running (Executando)}, o processo está efetivamente utilizando a CPU naquele instante.

No \texttt{top}, ambos os estados aparecem com o código \texttt{R}, pois o Linux agrupa processos prontos e em execução sob a mesma classificação. A alternância constante na ordem dos processos reflete as trocas de contexto realizadas pelo escalonador: a cada fatia de tempo, um processo diferente pode estar de fato na CPU, enquanto os demais aguardam na fila de prontos.

\subsection*{2. Ao executar pstree, você notou que o xeyes é filho de qual processo?}

O \texttt{xeyes} é filho do processo \texttt{bash}, que é o shell a partir do qual o comando foi executado. A hierarquia típica observada é:

\begin{lstlisting}
systemd -> terminal -> bash -> xeyes
\end{lstlisting}

Isso ocorre porque, no Linux, todo processo é criado por outro processo (seu pai) através da chamada de sistema \texttt{fork()}, formando uma árvore hierárquica.

\subsection*{3. Explique por que, ao aumentarmos o valor de nice para 15, o processo tende a usar menos CPU quando há outros processos ativos.}

O escalonador CFS (\textit{Completely Fair Scheduler}) do Linux utiliza o valor de \textit{nice} para calcular o peso relativo de cada processo na distribuição de tempo de CPU. O valor de \textit{nice} varia de $-20$ (maior prioridade) a $+19$ (menor prioridade), sendo 0 o padrão.

Com nice = 15, o processo recebe um peso significativamente menor em comparação com processos de nice = 0. Quando há outros processos ativos competindo pela CPU, o escalonador atribui fatias de tempo proporcionais ao peso de cada processo, fazendo com que o processo com nice 15 receba consideravelmente menos tempo de CPU.

Vale notar que, em um sistema ocioso (sem concorrência), mesmo um processo com nice alto pode utilizar toda a CPU disponível, pois não há outros processos disputando o recurso.

\subsection*{4. O que acontece com o registro de um processo no SO quando enviamos um sinal kill -9?}

O sinal \texttt{SIGKILL} (9) é tratado diretamente pelo kernel e não pode ser interceptado, bloqueado ou ignorado pelo processo. Ao recebê-lo, o kernel encerra o processo imediatamente, sem permitir a execução de \textit{handlers} de limpeza. Em seguida, o sistema operacional:

\begin{enumerate}
  \item Libera todos os recursos alocados ao processo (memória, arquivos abertos, conexões de rede);
  \item Notifica o processo pai através do sinal \texttt{SIGCHLD};
  \item Remove a entrada do processo na tabela de processos (PCB).
\end{enumerate}

Se o processo pai não executar a chamada \texttt{wait()} para recolher o código de saída do filho, o processo encerrado permanece temporariamente no estado \textit{zumbi} (\texttt{Z}) até ser recolhido.

\end{document}